\section*{Hoofdstuk \ref{ch:b3:orogenic-wedges} --- Overschuivingsgordels en orogene wiggen}

\begin{solution}{exo:b3:orogenic-wedges:1}
$\sin\ang{30} = 0.5$: $K = 2\times0.5/0.5 = 2$ en $\sigma'_1/\sigma'_3 = 1 +
K = 3$.
\end{solution}

\begin{solution}{exo:b3:orogenic-wedges:2}
$\phi = \arctan0.6 = \ang{31.0}$, $\sin\phi = 0.514$, $K = 2.12$. Op een
horizontale basis, $\alpha = 0.3/3.12 = \qty{0.0962}{rad} = \ang{5.5}$. Op een
hellende basis $\ang{3} = \qty{0.0524}{rad}$, $\alpha + \beta = 0.352/3.12 =
\qty{0.113}{rad} = \ang{6.5}$, dus $\alpha = \ang{3.5}$.
\end{solution}

\begin{solution}{exo:b3:orogenic-wedges:3}
Diepte $= (T - 10)/30$: apatiet (U--Th)/He 1.3 tot \qty{2.3}{km};
apatiet-fissiesporen 2.7 tot \qty{3.7}{km}; zirkoon-fissiesporen $270/30 = \qty{9.0}{km}$. Op
\qty{5}{km} diepte lag het gesteente op $10 + 150 = \qty{160}{\celsius}$: zirkoon
(U--Th)/He, die sluit bij 160 tot \qty{200}{\celsius}, 5.0 tot
\qty{6.3}{km} diep.
\end{solution}

\begin{solution}{exo:b3:orogenic-wedges:4}
Horizontale lengte $2/\tan\ang{30} = \qty{3.46}{km}$; langs de helling
$2/\sin\ang{30} = \qty{4.0}{km}$. De laag wordt opgetild door de rampshoogte,
\qty{2}{km}. De afsnede moet de hele ramp afleggen: een verplaatsing van
\qty{4.0}{km}.
\end{solution}

\begin{solution}{exo:b3:orogenic-wedges:5}
$\phi = \arctan1.03 = \ang{45.8}$, $K = 5.08$; $\beta = \qty{0.105}{rad}$.
Voor $\lambda = 0.4$: $\alpha + \beta = (0.105 + 0.85\times0.6)/(1 +
5.08\times0.6) = 0.615/4.05 = \qty{0.152}{rad} = \ang{8.7}$, $\alpha =
\ang{2.7}$. Voor $\lambda = 0.9$: $0.190/1.51 = \qty{0.126}{rad} = \ang{7.2}$,
$\alpha = \ang{1.2}$. Verhoging van de vloeistofdruk verlaagt de \omterm{def:b3:orogenic-wedges:critical-taper}{kritische wighoek}:
de wig wordt superkritisch, schuift op zijn basis naar voren zonder
te verdikken, en de voorkant rukt op tot accretie en erosie hem
weer naar de nieuwe wighoek brengen.
\end{solution}

\begin{solution}{exo:b3:orogenic-wedges:6}
$S = A/r = 6.0/1.5 = \qty{4.0}{km}$. Gelezen als \omterm{def:b3:orogenic-wedges:thrust-folds}{onthechtingsplooi}, $S = A/h =
6.0/3.0 = \qty{2.0}{km}$: de helft van de juiste waarde. De twee lezingen moeten
onderscheiden worden aan de vorm van de plooi en de aanwezigheid van een ramp.
\end{solution}

\begin{solution}{exo:b3:orogenic-wedges:7}
$K = 2.12$ en $1 + K(1 - 0.4) = 2.27$. De vergelijking geeft $\mu_b(1 -
\lambda_b) = (\alpha + \beta)\times2.27 - \beta$, met $\beta =
\qty{0.0349}{rad}$. Jura: $0.0611\times2.27 - 0.0349 = 0.10$. Zuid:
$0.218\times2.27 - 0.0349 = 0.46$. Byerlee bij hydrostatische druk:
$0.85\times0.6 = 0.51$. De basis in het zuiden ligt dicht bij die van Byerlee; de
Jurabasis is bijna vijf keer zwakker, zoals verwacht van een décollement in
haliet, waarvan de wrijving laag is, onder vloeistoffen zonder overdruk.
\end{solution}

\begin{solution}{exo:b3:orogenic-wedges:8}
$\kappa = \qty{e-6}{m^2/s} = \qty{31.6}{km^2/Myr}$, en $v =
\qty{1}{km/Myr}$: $\mathrm{Pe} = 30/31.6 = 0.951$. De factor $0.951/(1 -
e^{-0.951}) = 1.55$ verhoogt de gradiënt van 30 naar \qty{46.5}{\celsius/km}.
Voor \qty{65}{\celsius}: $1 - e^{-\mathrm{Pe}\,z/L} = (55/900)(1 - e^{-0.951})
= 0.0375$, dus $\mathrm{Pe}\,z/L = 0.0382$ en $z = \qty{1.21}{km}$, tegenover
$55/30 = \qty{1.83}{km}$ zonder erosie.
\end{solution}

\begin{solution}{exo:b3:orogenic-wedges:9}
$\mathrm dh/\mathrm dt = \qty{0.8}{km}/\qty{1.2}{Myr} = \qty{0.67}{mm/yr}$.
Met $f = 0.5$, $v = 0.67\times(1 - 0.5) = \qty{0.33}{mm/yr}$.
\end{solution}

\begin{solution}{exo:b3:orogenic-wedges:10}
Met $\alpha = 0$ en $u = 1 - \lambda$: $\beta(1 + Ku) = \beta + \mu_b u$,
dus $\beta Ku = \mu_b u$ en $\beta^* = \mu_b/K$, ongeacht $u$. Voor Taiwan,
$0.85/5.08 = \qty{0.167}{rad} = \ang{9.6}$. Op een steilere basis is de kritische
oppervlaktehelling is negatief: elke wig met een vlak of naar het voorland hellend
oppervlak is al superkritisch, en hij glijdt op zijn basis zonder te verdikken. Een
overschuivingsgordel kan daar zijn normale vorm niet aannemen.
\end{solution}

\begin{solution}{exo:b3:orogenic-wedges:11}
Als $\Lambda \to \infty$ en $v \to 0$, $m \to 0$, dus $e^{-mz_c} \to 1$, en
de geotherm is rechtlijnig, $G_0 = G(z_c) = G$: $f \to (G - \Gamma)/G = 1 -
\Gamma/G = 1 - 6.5/30 = 0.78$. De helling is dan $v/(1 - f) = vG/\Gamma =
4.6v$. Voor $v = \qty{1}{mm/yr}$, $v/2\kappa = \qty{0.0158}{km^{-1}}$. Op
\qty{5}{km}: $m = 0.0158 + \sqrt{0.0158^2 + 1.257^2} = \qty{1.273}{km^{-1}}$,
$e^{-1.273\times1.21} = 0.215$, $f = 40.0\times0.215/44.7 = 0.19$. Op
\qty{20}{km}: $m = 0.0158 + 0.3146 = \qty{0.330}{km^{-1}}$, $e^{-0.399} =
0.671$, $f = 40.0\times0.671/44.7 = 0.60$.
\end{solution}

\begin{solution}{exo:b3:orogenic-wedges:12}
De doorsnede is een driehoek met basis $W$ en hoogte $W\tan(\alpha + \beta)$:
$A = W^2\tan(\alpha + \beta)/2$. Met $\mathrm dA/\mathrm dt = F_a$ en $A =
0$ bij $t = 0$, $A = F_at$, wat $W$ geeft. Met $\tan\ang{6} = 0.105$:
$W = \sqrt{2\times40\times2/0.105} = \qty{39}{km}$ na \qty{2}{Myr}; $t =
60^2\times0.105/(2\times40) = \qty{4.7}{Myr}$ om \qty{60}{km} te bereiken.
\end{solution}

\begin{solution}{pb:b3:orogenic-wedges:1}
\textbf{1.} $4.5 - 1.5 = \qty{3.0}{km}$ hoog; $3.0/\tan\ang{25} =
\qty{6.43}{km}$ in horizontale lengte; $3.0/\sin\ang{25} = \qty{7.10}{km}$ langs
de helling.

\textbf{2.} Door de rampshoogte, \qty{3.0}{km}.

\textbf{3.} $S = A/r = 16.5/3.0 = \qty{5.5}{km}$.

\textbf{4.} Hetzelfde, \qty{16.5}{km^2}: in een \omterm{def:b3:orogenic-wedges:thrust-folds}{breukbuigplooi} is de overtollige oppervlakte
van elke laag boven de bovenste flat is $Sr$. Als \omterm{def:b3:orogenic-wedges:thrust-folds}{onthechtingsplooi}, $16.5/4.0 =
\qty{4.1}{km}$: een kwart te weinig.

\textbf{5.} $4\times(6.0 - 2.5) = \qty{14.0}{km}$, tweeënhalf keer de
\qty{5.5}{km} verschuiving op de frontale ramp.

\textbf{6.} De frontale schub, die het dichtst bij het voorland ligt: elke nieuwe ramp
vormt zich in de liggende wand vóór de vorige. Elke schub, langs zijn ramp omhooggedragen,
tilt de dakoverschuiving en de bedekking erboven op, net zoals een \omterm{def:b3:orogenic-wedges:thrust-folds}{breukbuigplooi}
zijn hangende wand optilt, zodat het dak over de stapel wordt geplooid.

\textbf{7.} $\alpha = \arctan(2.1/60) = \ang{2.00}$; $\beta =
\arctan(4.2/60) = \ang{4.00}$; wighoek \ang{6.0}.

\textbf{8.} $2.4 + 8.2 = \qty{10.6}{km}$, minder dan ongeveer \qty{15}{km}: de
basis is nog wrijvend, en de theorie van Coulomb is van toepassing.

\textbf{9.} $\phi = \ang{45.8}$, $\sin\phi = 0.717$, $K = 5.08$. Droog:
$\alpha + \beta = (0.0699 + 0.85)/6.08 = \qty{0.151}{rad} = \ang{8.7}$,
$\alpha = \ang{4.7}$.

\textbf{10.} $\lambda = 1000/2500 = 0.40$: $(0.0699 + 0.51)/4.05 =
\qty{0.143}{rad} = \ang{8.2}$, $\alpha = \ang{4.2}$. De waargenomen \ang{2.0}
zou de wig sterk subkritisch maken, zodat die eerder verdikt dan oprukt:
de fluïda moeten overdruk hebben.

\textbf{11.} Met $T = \qty{0.1049}{rad}$ en $u = 1 - \lambda$, $T(1 + Ku) =
\beta + \mu_bu$ geeft $u = (\beta - T)/(TK - \mu_b) = (0.0699 -
0.1049)/(0.5327 - 0.85) = 0.110$: $\lambda = 0.89$.

\textbf{12.} $\sigma_v = 2500\times9.81\times10\,600 = \qty{260}{MPa}$; $P =
0.89\times260 = \qty{231}{MPa}$; hydrostatisch $1000\times9.81\times10\,600 =
\qty{104}{MPa}$. De fluïda dragen \qty{127}{MPa} meer dan een kolom
water zou dragen.

\textbf{13.} $(0.0699 + 0.85\times0.3)/(1 + 5.08\times0.3) = 0.325/2.52 =
\qty{0.129}{rad} = \ang{7.4}$, $\alpha = \ang{3.4}$. De wig, bij \ang{2.0},
is nu subkritisch: hij houdt op met oprukken en verdikt door
\omterm{def:b3:orogenic-wedges:thrust-geometry}{overschuiving buiten de normale volgorde} achter zijn voorkant, totdat zijn oppervlaktehelling \ang{3.4} bereikt.

\textbf{14.} $F_a = \qty{10}{km/Myr}\times\qty{4.0}{km} = \qty{40}{km^2/Myr}$,
dat is \qty{4.0e4}{m^3/yr} per kilometer van de gordel.

\textbf{15.} $\bar e = 40/60 = \qty{0.67}{km/Myr} = \qty{0.67}{mm/yr}$.

\textbf{16.} $40\times100 = \qty{4000}{km^3/Myr} = \qty{4.0e6}{m^3/yr}$, of
$\qty{1.0e10}{kg/yr}$ (10 miljoen ton per jaar).

\textbf{17.} $\mathrm dW/\mathrm dt = F_a/(W\tan(\alpha + \beta)) =
40/(60\times0.105) = \qty{6.3}{km/Myr}$.

\textbf{18.} Nee: in stationaire toestand moet de erosie nog steeds de hele
inkomende flux verwijderen, dus de gemiddelde snelheid blijft \qty{0.67}{mm/yr}. Maar het gesteente dat
onder de achterkant wordt toegevoegd, wordt daar opgetild: erosie en exhumatie zouden het snelst
optreden boven de achterkant, waar de jongste afkoelingsouderdommen zouden worden gevonden.

\textbf{19.} $v = \qty{0.667}{km/Myr}$: $\mathrm{Pe} = 0.667\times30/31.6 =
0.634$; $G_0 = 30\times0.634/(1 - e^{-0.634}) = 30\times1.35 =
\qty{40.5}{\celsius/km}$.

\textbf{20.} $1 - e^{-\mathrm{Pe}\,z/L} = (55/900)(1 - e^{-0.634}) = 0.0287$,
dus $\mathrm{Pe}\,z/L = 0.0291$ en $z_c = \qty{1.38}{km}$; $G(z_c) =
40.5\,e^{-0.0291} = \qty{39.3}{\celsius/km}$.

\textbf{21.} $1.38/0.667 = \qty{2.07}{Myr}$, de ouderdom van het monster op
\qty{2.00}{km}: op de gemiddelde hoogte is de doorbuiging van de isotherm nul.

\textbf{22.} $v/2\kappa = 0.667/63.1 = \qty{0.0106}{km^{-1}}$, $2\pi/\Lambda =
\qty{0.314}{km^{-1}}$: $m = 0.0106 + 0.3143 = \qty{0.325}{km^{-1}}$;
$e^{-0.325\times1.38} = 0.639$; $f = (40.5 - 6.5)\times0.639/39.3 = 0.552$.

\textbf{23.} De gemiddelde hoogte is \qty{2.00}{km} en de gemiddelde ouderdom
\qty{2.068}{Myr}; $\sum\Delta t\,\Delta h = 0.823$ en $\sum\Delta t^2 =
0.552$, dus de helling is \qty{1.49}{km/Myr}: een schijnbare snelheid van
\qty{1.49}{mm/yr}.

\textbf{24.} $v = 1.49\times(1 - 0.552) = \qty{0.667}{mm/yr}$.

\textbf{25.} De \omterm{def:b3:orogenic-wedges:thermochronology}{thermochronometer} en de fluxbalans komen overeen: de wig is
in stationaire toestand, en de schijnbare helling voor waar aannemen zou
de snelheid meer dan twee keer te hoog hebben ingeschat. De exhumatiesnelheid van de wig Kestrel
is \textbf{\qty[text-series-to-math=true]{0.67}{mm/yr}}.
\end{solution}
